\documentclass[10pt]{article}
\usepackage[utf8]{inputenc}
\usepackage[T1]{fontenc}
\usepackage{amsmath, amssymb}
\usepackage{hyperref}
\usepackage[margin=1in]{geometry}
\usepackage{xcolor}
\usepackage{tcolorbox}
\usepackage{fancyhdr}
\usepackage{titlesec}
\usepackage{longtable}
\usepackage{parskip}

\definecolor{gapred}{RGB}{178,24,43}
\definecolor{gaporange}{RGB}{217,95,2}
\definecolor{codebg}{RGB}{245,245,245}

\newtcolorbox{gapbanner}{
  colback=gapred!6, colframe=gapred, boxrule=1pt, arc=2pt,
  fonttitle=\bfseries, title=READ THIS FIRST}
\newcommand{\gap}[1]{\textcolor{gapred}{\textbf{[GAP]} #1}}
\newcommand{\flag}[1]{\textcolor{gaporange}{\textbf{[FLAG]} #1}}

\newcommand{\rdtopic}[1]{\clearpage\section{#1}}
\newcommand{\rdsub}[1]{\subsection*{#1}}
\newenvironment{rdusage}{\begin{quote}\ttfamily\small}{\end{quote}}
\newenvironment{rdcode}{\begin{quote}\ttfamily\small}{\end{quote}}

\hypersetup{colorlinks=true, linkcolor=blue!50!black, urlcolor=blue!50!black}

\pagestyle{fancy}
\fancyhf{}
\fancyhead[L]{BayGMST}
\fancyhead[R]{\thepage}
\renewcommand{\headrulewidth}{0.4pt}

\title{Package `BayGMST'}
\date{Draft reference manual --- \today}
\author{}

\begin{document}
\maketitle
\thispagestyle{empty}

\begin{center}
\begin{tabular}{ll}
\textbf{Type:} & Package \\
\textbf{Title:} & Bayesian Reconstruction of Global Mean Surface Temperature \\
\textbf{Version:} & 0.1.0 \\
\textbf{Maintainer:} & Tyler Bagwell \texttt{<teb6@rice.edu>} \\
\textbf{License:} & GPL ($\geq$ 3) \\
\end{tabular}
\end{center}

\vspace{1em}
\begin{gapbanner}
This is a \textbf{hand-typeset preview}, not the real, authoritative CRAN
reference manual. It was produced in a sandbox with \textbf{no R
installation available}, so it could not be generated the normal way
(\texttt{devtools::document()} to build \texttt{man/*.Rd}, then
\texttt{devtools::build\_manual()} / \texttt{R CMD Rd2pdf} to render this
PDF from those \texttt{.Rd} files). Content below was written by hand to
mirror that output as closely as possible, from the same source
descriptions as the \texttt{roxygen2} comments in \texttt{R/}.

\textbf{Before treating this as final:} run the two commands above (see
\texttt{CRAN-READINESS.md} at the package root for the full checklist) and
let the real output supersede this file.

Text marked \gap{like this} is a genuine information gap I could not
resolve myself. Text marked \flag{like this} is a decision, judgment call,
or unresolved issue from the original code that needs Tyler's or Julien's
sign-off, not something I decided unilaterally.
\end{gapbanner}

\clearpage
\tableofcontents

\rdtopic{R topics documented:}
\begin{quote}
\begin{tabular}{ll}
BayGMST-package & BayGMST: Bayesian Reconstruction of Global Mean Surface Temperature \\
cv\_baygmst & K-fold cross-validate the BayGMST model over the instrumental period \\
fit\_baygmst & Fit the BayGMST Bayesian hierarchical AR(1) reconstruction model \\
plot\_cv & Plot k-fold cross-validation reconstructions \\
plot\_posterior\_densities & Posterior density plots for a fitted BayGMST model \\
plot\_reconstruction & Plot a BayGMST reconstruction \\
plot\_trace & Trace plots for a fitted BayGMST model \\
print.baygmst\_cv & Print method for \texttt{"baygmst\_cv"} objects \\
print.baygmst\_fit & Print a fitted BayGMST model \\
print.summary.baygmst\_fit & Print method for \texttt{"summary.baygmst\_fit"} objects \\
reconstruct & Extract a tidy reconstruction from a fitted BayGMST model \\
reduce\_proxies & Reduce a multi-proxy network to a single representative proxy series \\
summary.baygmst\_fit & Summarize a fitted BayGMST model \\
transform\_forcings & Transform raw radiative forcings for use with \texttt{fit\_baygmst()} \\
\end{tabular}
\end{quote}

% ============================================================
\rdtopic{BayGMST-package}
\textbf{BayGMST: Bayesian Reconstruction of Global Mean Surface Temperature}

\rdsub{Description}
Fits a Bayesian hierarchical AR(1) state-space model relating global mean
surface temperature (GMST) to a reduced paleoclimate proxy network,
instrumental temperature observations, and radiative forcings (greenhouse
gases, volcanic aerosols, and solar irradiance). Missing pre-instrumental
temperatures are estimated jointly with the model parameters using Stan via
the \texttt{instantiate} and \texttt{cmdstanr} packages.

\rdsub{Model fitting requires CmdStan}
\texttt{fit\_baygmst()} and \texttt{cv\_baygmst()} compile and sample a Stan
model, which requires a working CmdStan installation. Check with
\texttt{instantiate::stan\_cmdstan\_exists()} first. \texttt{reduce\_proxies()}
and \texttt{transform\_forcings()} do not require CmdStan.

\rdsub{Provenance}
Restructured from a collection of top-level analysis scripts (see
\texttt{NEWS.md} for exactly what changed, and
\texttt{inst/legacy-scripts/} for the untouched originals).
\gap{The \texttt{Description:} field's exact public-facing wording in
\texttt{DESCRIPTION} was drafted by me from the README and the thesis
proposal's Project II motivation section --- it has not been reviewed or
approved by Tyler or Julien. Please read and edit it before any submission.}

\rdsub{Author(s)}
Tyler Bagwell \texttt{<teb6@rice.edu>} (maintainer), Julien Emile-Geay
\texttt{<julieneg@usc.edu>}.
\flag{These are the only two contributors visible in \texttt{git log}; add
others if appropriate.}

% ============================================================
\rdtopic{transform\_forcings}
\textbf{Transform raw radiative forcings for use with \texttt{fit\_baygmst()}}

\rdsub{Description}
Applies the forcing transformations used throughout the BayGMST pipeline:
volcanic forcing is mapped to a negative saturating form, greenhouse-gas
(CO2) forcing is log-transformed relative to a reference concentration, and
solar forcing is centered on its own mean. Ported from the normalization
block near the top of the original \texttt{BayGMST\_v1.0.R} script.

\rdsub{Usage}
\begin{rdusage}
transform\_forcings(G, V, S, co2\_c0 = 280, co2\_coef = 5.35, vol\_coef = 25)
\end{rdusage}

\rdsub{Arguments}
\begin{quote}
\begin{tabular}{@{}p{3cm}p{11cm}@{}}
\texttt{G} & Numeric vector of raw greenhouse-gas forcing (e.g.\ atmospheric CO2 concentration in ppm). \\[0.3em]
\texttt{V} & Numeric vector of raw volcanic forcing (e.g.\ aerosol optical depth). Same length as \texttt{G}. \\[0.3em]
\texttt{S} & Numeric vector of raw solar forcing (e.g.\ total solar irradiance). Same length as \texttt{G}. \\[0.3em]
\texttt{co2\_c0} & Reference/baseline CO2 concentration used in the log transform. Default 280 (ppm). \\[0.3em]
\texttt{co2\_coef} & Multiplier applied to the log CO2 ratio. Default 5.35. \\[0.3em]
\texttt{vol\_coef} & Multiplier applied to the volcanic saturating transform. Default 25. \\
\end{tabular}
\end{quote}

\rdsub{Details}
\[ V' = -|\texttt{vol\_coef}| \times (1 - e^{-V}), \qquad
   G' = \texttt{co2\_coef} \times \log(G / \texttt{co2\_c0}), \qquad
   S' = S - \bar{S} \]

\rdsub{Value}
A list with three numeric vectors, \texttt{G}, \texttt{V}, and \texttt{S} --
the transformed forcings.

\rdsub{Examples}
\begin{rdcode}
fc <- transform\_forcings(\\
\ \ G = c(280, 300, 340, 400),\\
\ \ V = c(0, 0.05, 0.02, 0.10),\\
\ \ S = c(1360, 1361, 1360.5, 1361.2)\\
)\\
fc\$G
\end{rdcode}

% ============================================================
\rdtopic{reduce\_proxies}
\textbf{Reduce a multi-proxy network to a single representative proxy series}

\rdsub{Description}
Collapses a matrix of paleoclimate proxy records into a "representative
proxy" series suitable for \texttt{fit\_baygmst()}, following the
segmented-composite approach of Barboza et al.\ (2019), ported from
\texttt{utils/PAGES2k\_reducedProxy\_UNSC.R}. The proxy network is broken
into overlapping time segments of \texttt{chunk} years; within each
segment, available proxies are regressed against instrumental temperature
during \texttt{calib\_years} using the chosen \texttt{method}. The final
composite series averages across all segments covering a given year.

\rdsub{Usage}
\begin{rdusage}
reduce\_proxies(proxy\_matrix, years, temp\_calib, calib\_years,\\
\ \ method = c("PCR", "LASSO", "sPCR", "SPLS", "SIR"),\\
\ \ chunk = 250, target\_adj\_r2 = 0.7, max\_na\_frac = 0.05,\\
\ \ n\_components = 3, spls\_eta = seq(0.1, 0.9, 0.1))
\end{rdusage}

\rdsub{Arguments}
\begin{quote}
\begin{tabular}{@{}p{3.2cm}p{10.8cm}@{}}
\texttt{proxy\_matrix} & Numeric matrix/data.frame, one row per year (matching \texttt{years}), one column per proxy. No year column. \\[0.3em]
\texttt{years} & Integer vector of calendar years, one per row of \texttt{proxy\_matrix}. \\[0.3em]
\texttt{temp\_calib} & Numeric vector, instrumental temperature during \texttt{calib\_years}. \\[0.3em]
\texttt{calib\_years} & Integer vector, calibration years (subset of \texttt{years}; \texttt{length(calib\_years) == length(temp\_calib)}). \\[0.3em]
\texttt{method} & One of \texttt{"PCR"} (default), \texttt{"LASSO"}, \texttt{"sPCR"}, \texttt{"SPLS"}, or \texttt{"SIR"}. \textbf{\texttt{"SIR"} is experimental} --- see Details. \\[0.3em]
\texttt{chunk} & Segment length in years. Default 250. Every segment must fully contain \texttt{calib\_years}, or \texttt{reduce\_proxies()} errors (see Details). \\[0.3em]
\texttt{target\_adj\_r2} & For \texttt{"PCR"}: adjusted-$R^2$ stopping threshold for adding principal components. Default 0.70. \\[0.3em]
\texttt{max\_na\_frac} & Proxies missing more than this fraction of a segment are dropped from it. Default 0.05. \\[0.3em]
\texttt{n\_components} & For \texttt{"sPCR"} only: number of supervised components (scalar or per-segment vector). Default 3. \\[0.3em]
\texttt{spls\_eta} & For \texttt{"SPLS"} only: sparsity grid passed to \texttt{spls::cv.spls()}. \\
\end{tabular}
\end{quote}

\rdsub{Details}
\textbf{Method support.} \texttt{"PCR"}, \texttt{"LASSO"}, and \texttt{"SPLS"}
are fully adaptive and supported. \texttt{"sPCR"} is supported but requires
\texttt{n\_components} to be chosen by the caller (not cross-validated,
matching the original script).

\flag{\texttt{"SIR"} is a simplified, \emph{unvalidated} reimplementation
using a single sliced-inverse-regression direction from \texttt{dr}, with
no per-segment variable pre-selection. It intentionally does \emph{not}
reproduce the original script's SIR branch, which called an undocumented
\texttt{edrSelec()} step from the non-CRAN package `edrGraphicalTools' that
the original author flagged as never fully working ("NEED TO FIGURE OUT HOW
TO INSTALL edrGraphicalTools!!"). \texttt{config.yml} in the original
repository likewise documented SIR as "still under construction." Treat
\texttt{method = "SIR"} output with corresponding caution.}

\flag{The segmentation/calibration-window constraint (\texttt{chunk} must
be small enough that even the shortest segment still spans all of
\texttt{calib\_years}) mirrors an explicit \texttt{stop()} check in the
original script. I ported that check explicitly rather than letting a
mismatched-length input fail with a cryptic error deep inside \texttt{lm()}
--- but the underlying design constraint itself is inherited from the
original code, not something I introduced.}

Required packages: \texttt{"LASSO"} needs \texttt{glmnet}; \texttt{"sPCR"}
needs \texttt{superpc}; \texttt{"SPLS"} needs \texttt{spls}; \texttt{"SIR"}
needs \texttt{dr}. All are \texttt{Suggests}, only required for the method
actually called.

\rdsub{Value}
A list with \texttt{segments} (data.frame, one column per segment),
\texttt{composite} (data.frame with columns \texttt{year} and \texttt{RP1},
the series to pass to \texttt{fit\_baygmst()}), and \texttt{method}.

% ============================================================
\rdtopic{fit\_baygmst}
\textbf{Fit the BayGMST Bayesian hierarchical AR(1) reconstruction model}

\rdsub{Description}
Prepares the data list expected by the \texttt{baygmst} Stan model and
samples from it with \texttt{cmdstanr}, via \texttt{instantiate}. Ported
from the data-preparation and \texttt{mod\$sample()} block of the original
\texttt{BayGMST\_v1.0.R} script, restructured to take explicit arguments
instead of reading \texttt{config.yml} and hardcoded relative paths, and to
return an object rather than leaving variables in the global environment.

\rdsub{Usage}
\begin{rdusage}
fit\_baygmst(proxy, instrumental\_T, forcing\_G, forcing\_V, forcing\_S, years,\\
\ \ chains = 4, parallel\_chains = 2, iter\_warmup = 500,\\
\ \ iter\_sampling = 1500, seed = NULL, ...)
\end{rdusage}

\rdsub{Arguments}
\begin{quote}
\begin{tabular}{@{}p{3.4cm}p{10.6cm}@{}}
\texttt{proxy} & Representative proxy series -- see \texttt{reduce\_proxies()}. No missing values. \\[0.3em]
\texttt{instrumental\_T} & Instrumental temperature, \texttt{NA} for pre-instrumental years to reconstruct. \\[0.3em]
\texttt{forcing\_G, forcing\_V, forcing\_S} & Already-transformed forcings -- see \texttt{transform\_forcings()}. No missing values. \\[0.3em]
\texttt{years} & Calendar years, defining row order for all arguments above. \\[0.3em]
\texttt{chains, parallel\_chains, iter\_warmup, iter\_sampling, seed} & Passed to \texttt{cmdstanr}'s \texttt{\$sample()}. \texttt{iter\_sampling} $\geq$ 1000 required. \\[0.3em]
\texttt{...} & Additional arguments passed to \texttt{\$sample()}. \\
\end{tabular}
\end{quote}

\rdsub{Value}
An object of class \texttt{"baygmst\_fit"}: a list with \texttt{fit} (the
\texttt{CmdStanMCMC} object), \texttt{years}, \texttt{idx\_obs},
\texttt{idx\_mis}, \texttt{data}, and \texttt{call}.

\rdsub{Note}
\gap{This function does not write any files by default (no auto-\texttt{ggsave()},
no auto-CSV export) -- a deliberate deviation from the original script,
which wrote directly to \texttt{cfg\$folder\_paths\$...}. This matches CRAN
policy on functions not writing to the filesystem by default, but it does
mean any code that relied on those side-effect files (e.g.\
\texttt{outputs/reconstructions/fit\_post\_summaries.csv}) needs to be
updated to call \texttt{write.csv()} itself on the returned object.}

% ============================================================
\rdtopic{cv\_baygmst}
\textbf{K-fold cross-validate the BayGMST model over the instrumental period}

\rdsub{Description}
Partitions the instrumental years of \texttt{instrumental\_T} into
\texttt{nfold} contiguous blocks; for each, refits the model with that
block hidden, and reports recovery accuracy. Ported from the fold loop in
\texttt{cv\_v0.1.R}.

\rdsub{Usage}
\begin{rdusage}
cv\_baygmst(proxy, instrumental\_T, forcing\_G, forcing\_V, forcing\_S, years,\\
\ \ nfold = 3, chains = 4, parallel\_chains = 2, iter\_warmup = 500,\\
\ \ iter\_sampling = 1500, seed = NULL, ...)
\end{rdusage}

\rdsub{Arguments}
Same as \texttt{fit\_baygmst()}, plus \texttt{nfold} (number of folds,
default 3).

\rdsub{Value}
An object of class \texttt{"baygmst\_cv"}: \texttt{folds} (combined
data.frame of held-out predictions), \texttt{r2} and \texttt{mse}
(per-fold, posterior mean of the Stan-computed \texttt{r2\_cv}/\texttt{mse\_cv}
generated quantities), \texttt{years}, \texttt{call}.

\rdsub{Note}
\flag{Two blocks in the original \texttt{cv\_v0.1.R} were dropped, not
ported: (1) a \texttt{Forcings.projections}/RCP-scenario block that was
built but never referenced again (dead code), and which additionally
hardcoded a path to a specific contributor's machine
(\texttt{/Users/tylerbagwell/...}); and (2) a large commented "OLD"
duplicate of the fold loop at the end of the file. Confirm nothing
downstream still depended on either.}

\flag{This implementation reports the fully Bayesian posterior mean of the
Stan model's own \texttt{r2\_cv}/\texttt{mse\_cv} generated quantities. The
original (non-dead) R loop instead computed a plug-in point-estimate
$R^2$ in R from the posterior mean reconstruction. The two will be close
but not identical -- flagging the substitution explicitly rather than
silently changing what "R2" means.}

% ============================================================
\rdtopic{reconstruct}
\textbf{Extract a tidy reconstruction from a fitted BayGMST model}

\rdsub{Description}
Builds a tidy data.frame combining observed instrumental temperatures with
the posterior reconstruction of pre-instrumental years -- equivalent to the
\texttt{gmst\_reconstruction\_data.csv} the original script wrote to disk as
a side effect, but returned as an object instead.

\rdsub{Usage}
\begin{rdusage}
reconstruct(fit, probs\_inner = c(0.16, 0.84), probs\_outer = c(0.025, 0.975))
\end{rdusage}

\rdsub{Value}
A data.frame with columns \texttt{year}, \texttt{type}
(\texttt{"instrumental"}/\texttt{"reconstruction"}), \texttt{T\_obs},
\texttt{T\_mean}, \texttt{T\_lo\_inner}, \texttt{T\_hi\_inner},
\texttt{T\_lo\_outer}, \texttt{T\_hi\_outer}.

% ============================================================
\rdtopic{plot\_reconstruction, plot\_trace, plot\_posterior\_densities, plot\_cv}
\textbf{Plotting functions}

\rdsub{Description}
\begin{itemize}
\item \texttt{plot\_reconstruction(fit, title = "GMST Reconstruction")} --
  ported from the \texttt{p\_ts} block of \texttt{BayGMST\_v1.0.R}.
\item \texttt{plot\_trace(fit, parameters = ...)} -- ported from the
  trace-plot block.
\item \texttt{plot\_posterior\_densities(fit, parameters = ...)} -- ported
  from the posterior-histogram block.
\item \texttt{plot\_cv(cv)} -- ported from the \texttt{p\_ts\_cv} block of
  \texttt{cv\_v0.1.R}.
\end{itemize}
All four return a \texttt{ggplot} object instead of calling
\texttt{ggsave()}/auto-printing as a side effect; save with
\texttt{ggplot2::ggsave()} if a file is wanted.
\flag{The original script combined the reconstruction and posterior-density
panels into one figure via \texttt{patchwork} (\texttt{p\_ts + p\_hist}).
Splitting these into separate functions means that combination is no
longer automatic -- users who want the original combined layout can do
\texttt{plot\_reconstruction(fit) | plot\_posterior\_densities(fit)}
themselves via \texttt{patchwork} (moved to \texttt{Suggests}, since the
package itself no longer calls it directly).}

% ============================================================
\rdtopic{print/summary methods}
\textbf{S3 methods}

\rdsub{Description}
\begin{itemize}
\item \texttt{print.baygmst\_fit(x, ...)} -- one-line-per-field overview
  (year range, chains/draws, call).
\item \texttt{summary.baygmst\_fit(object, ...)} -- posterior parameter
  summary (\texttt{alpha0, alpha1, phi\_R, phi\_T, beta0, betaG, betaS,
  betaV, sigma\_y, sigma\_z}) plus instrumental-period fit performance
  (RMSE, MAE, Bias, Correlation, $R^2$, detrended $R^2$), replacing the
  \texttt{perf\_stats}/\texttt{R2\_dt} block of \texttt{BayGMST\_v1.0.R}.
\item \texttt{print.summary.baygmst\_fit(x, ...)} -- pretty-prints the
  above.
\item \texttt{print.baygmst\_cv(x, ...)} -- per-fold and median $R^2$.
\end{itemize}

\clearpage
\section*{Known limitations index}
\gap{The unresolved possible bug in \texttt{src/stan/baygmst.stan}
(\texttt{y\_ins\_fitted} block, \texttt{t == 1} branch: \texttt{z[NT\_mis]}
used as an index where \texttt{z[idx\_mis[NT\_mis]]} looks more likely
correct) is carried forward \emph{unchanged} from the original code. This
restructuring does not alter model math without Tyler/Julien's statistical
sign-off. See \texttt{NEWS.md} for the exact location and reasoning.}

\flag{Several \texttt{library()} calls in the original scripts
(\texttt{car}, \texttt{fda}, \texttt{ggmap}, \texttt{maps}) had no evident
corresponding usage in the code as written and were not carried into
\texttt{DESCRIPTION}. If a hidden use is found, re-add them.}

\gap{\texttt{DESCRIPTION}'s \texttt{RoxygenNote} field is a placeholder
(\texttt{7.3.2}) -- \texttt{devtools::document()} will set the real value.}

\gap{\texttt{NAMESPACE} and every file in \texttt{man/} are hand-authored
placeholders (see the header comment in each), not actual
\texttt{roxygen2} output, because this restructuring was done without an R
installation available to run it. \texttt{CRAN-READINESS.md} has the exact
commands to regenerate and verify them for real.}

\end{document}
